% Created 2026-10-02 Fri 14:40
% Intended LaTeX compiler: xelatex
\documentclass[12pt]{article}
\usepackage{graphicx}
\usepackage{longtable}
\usepackage{wrapfig}
\usepackage{rotating}
\usepackage[normalem]{ulem}
\usepackage{amsmath}
\usepackage{amssymb}
\usepackage{capt-of}
\usepackage{hyperref}
\usepackage[a4paper,margin=1in]{geometry}
\usepackage{xurl}
\usepackage{graphicx}
\usepackage{parskip}
\usepackage{titling}
\makeatletter
\renewcommand{\maketitle}{%
\begin{center}
{\LARGE\@title\par}
\vspace{0.5em}
{\large\@author\par}
\vspace{0.5em}
{\large\@date\par}
\end{center}
}
\makeatother
\author{Filipa Calado}
\date{}
\title{Teaching Statement}
\hypersetup{
 pdfauthor={Filipa Calado},
 pdftitle={Teaching Statement},
 pdfkeywords={},
 pdfsubject={},
 pdfcreator={Emacs 29.3 (Org mode 9.6.15)}, 
 pdflang={English}}
\begin{document}

\maketitle
This teaching statement covers my teaching philosophy, the courses
I've taught at Pratt, mentoring and advising, and future plans.

\section*{Philosophy}
\label{sec:orgf0427c9}
As a language scholar who uses uses digital tools, I have experience
teaching in two general areas: English and computer science. In
English and related studies, I have taught courses on academic
writing, critical thinking, and Latinx literature. On the side of
computer science, I have taught courses Python programming, natural
language processing, data science, and machine learning. 

From this background, I have developed a teaching philosophy that
blends together humanist and computational skills, taking a humanist
approach to computational tools and objects. This philosophy can be
organized into three core principles:

\begin{quote}
\begin{enumerate}
\item the close is critical;
\item writing as an active learning method; and
\item knowledge is always contingent
\end{enumerate}
\end{quote}

In what follows, I explain these core principles in more detail.
\subsection*{1) the close is critical:}
\label{sec:org17f21b7}
In my courses, regardless of their format or focus, I teach close
reading as an essential critical practice. Originating from literary
studies, close reading is a particular kind of qualitative analysis
that sustains attention on a very small selection of evidence, such as
traditionally, a passage in a text. However, I find that close reading
helps students not only to analyze evidence, like textual passages,
depth, but also to generate new ways of using and interpreting all
kinds of data.

For example, in my INFO 664: Programming Cultural Heritage course, I
have students consider large datasets on an intimate scale,
demonstrating how attention to detail can reveal new material for
analysis and critique. In my INFO 664 class from last fall, I had
several students develop projects which used Small Language Models
(SLMs) to examine gender bias in popular culture,\footnote{Jace Steiner, "Considering Transformers Models as a Tool for
Character Analysis with Finn the Human." Fall 2025.
\url{https://xxsteinxx.github.io/chara-analysis\_at/}.} to study
discourses of toxic online spaces,\footnote{Ellen Connors. "Machine Divinity." Fall 2025.
\url{https://econn170.github.io/machine-divinity/index.html}.} and to surface patterns in
sensationalist language\footnote{Alana Maisel. "The New York Post's Environment and Wildlife."
Fall 2025.
\url{https://amaisel1.github.io/nypost-environment-and-wildlife/aboutproject.html}.}---three projects which ended up winning
the 2026 Info Show's "People's Choice Award."

\subsection*{2) writing as an active learning method:}
\label{sec:org21478f1}
I am a strong proponent of writing as a powerful and irreplaceable
skill for professional work. All of my courses emphasize the
development of this skill---all the more forcefully when emerging AI
technologies threaten to reduce students' intellectual capacities to
the capacities of any given tool.

Writing is, in my experience, the most efficient method of active
learning: it requires externalizing one's inner understanding in a
coherent and material way, while also exposing it to alternative
perspectives and real-world issues. In my courses, I teach writing
both through low stakes, daily formats like freewriting exercises, and
in long term assignments like formal papers. For example, in my INFO:
656 Machine Learning course, the major assignment is a conference
paper that students complete in groups and submit to a machine
learning (ML) conference. This assignment not only keeps them up to
date with recent advances in a rapidly changing field, it also
encourages students to think deeply about their contributions to ML
research within a professionalizing context.

\subsection*{3) knowledge is always contingent:}
\label{sec:org9c49250}
One of the first lessons I learned in the classroom is the importance
meeting students where they are, regardless of their previous
training. One way of doing so is to invite students to use their lived
experience as a ground for understanding and critique---to assure them
that their own histories, combined with reflection, constitute a
valuable kind of knowledge. A necessary corollary to knowledge as
experience is to accept that there can be no single perspective for
any question or issue. Rather than a single perspective, there are
degrees of visibility made available through habit and privilege. In
my teaching, I try to expand what is visible through the exposure to a
wide range of theoretical points of view, particularly in the
interdisciplinary selection of readings from my courses like INFO 601,
INFO 656, and INFO 657. Readings in these courses explore fields such
as Black Feminist Studies, Archival Studies, and AI Ethics research,
showing differences in initial questions, assumptions, approaches, and
their resulting conclusions.

\section*{Courses \& Course Evaluations}
\label{sec:orgd04f35b}
The following are a list of courses that I have taught since beginning
my position at Pratt in Fall 2024. Included for each course is
information about the semesters taught and number of students served,
links to course materials like syllabi and github repositories, and
excerpts from teaching evaluations.

\begin{itemize}
\item INFO 656: Machine Learning
\begin{itemize}
\item Semesters taught: Fall 2026
\item Students served: 11
\item Course github: \url{https://github.com/gofilipa/656}
\item Course syllabus:
\url{https://github.com/gofilipa/656/blob/main/syllabus.pdf}
\end{itemize}
\item INFO 657: Digital Humanities
\begin{itemize}
\item Semesters taught: Spring 2026
\item Students served:  11
\item Couse syllabus: \url{https://bit.ly/657-spring2026}
\item Excerpts from course evaluations::
\begin{itemize}
\item \emph{Professor Calado's welcoming approach to our discussions helped
to foster a mutually supportive classroom environment in which
we were all genuinely interested in each other's work, and
cheering each other on. I thought the guest speakers were
excellent. I loved the range of readings and the chance to get
introduced to so many DH tools.}
\item \emph{1. I loved the subject matter \& topics Dr. Calado choose
throughout the course, it really sparked my interest in DH and
helped me understand theoretical approaches to DH 2. I think the
assignments she chose also helped me understand approaching and
thinking about DH projects. If I were taking the advanced DH
course, I would feel very prepared for it with what I had
learned from this class. 3. I'm really happy she brough in
guest-speakers from the field, doing this also helped me
understand how I could possibly incorporate DH into my career
path and gave me more options to exlpore.}
\end{itemize}
\end{itemize}
\item INFO 697: Special Topics - How to Build a Bot
\begin{itemize}
\item Semesters taught: Summer 2025
\item Students served: 6
\item Course github: \url{https://github.com/gofilipa/botz}
\item Course syllabus: \url{https://github.com/gofilipa/botz/blob/main/syllabus.pdf}
\item Excerpts from course evaluations:
\begin{itemize}
\item \emph{Filipa is a great instructor and I've really enjoyed this
class, each week was super interesting and I really appreciate
her perspectives on coding, automation, and pretty much
everything!}
\item \emph{1. The conjunction of critical thinking supported by readings
with technical learning. 2. The hands-on aspect of the course
truly guided us in acquiring a powerful skill. 3. The open
dialogue and flexibility to adapt the program throughout the
entire time.}
\item \emph{I feel that now I have a superpower, and I'm eager to keep
practicing!}
\end{itemize}
\end{itemize}
\item INFO 664: Programming Cultural Heritage
\begin{itemize}
\item Semesters taught: Fall 2024, Spring 2025, Fall 2025, Spring 2026,
Fall 2026
\item Students served: 55
\item Course syllabus: \url{https://bit.ly/664-fall26}
\item Python notebooks: \url{https://gofilipa.github.io/664}
\item Student projects: \url{https://gofilipa.github.io/664-projects}
\item Excerpts from course evaluations:
\begin{itemize}
\item \emph{I think Prof. Calado structured the class really well and this
is the first teacher I had that has effectively taught me how to
code. She makes complex topics easy to understand no matter what
experience or background a student may or may not have with
coding.}
\item \emph{This course is really fun. Even if you aren't trying to be a
programmer, you learn really useful skills and it helps you
understand the internet and data on a deeper level. Prof. Calado
is very encouraging and very available. She's clearly in her
element when it comes to Python.}
\item \emph{I would make this a required course for all LIS students. I
wish I would have taken this course sooner, of all classes this
one has impacted my learning outcomes most.}
\item \emph{I appreciated the chance to openly discuss code/AI/programming
during class and one on one with the professor, learning about a
large breadth of uses for code in personal work/research, and
having a final where we could be creative and pursue a personal
interest.}
\end{itemize}
\end{itemize}
\item INFO 601: Foundations of Information 
\begin{itemize}
\item Semesters taught: Fall 2024, Spring 2025, Fall 2025, Spring 2026,
Fall 2026
\item Students served: 103
\item Course syllabus: \url{https://bit.ly/601-fall26}
\item Excerpts from course evaluations:
\begin{itemize}
\item \emph{Professor Calado’s Foundation of Information course was filled
with a lot of intimating topics and readings. However, Professor
Calado’s approach each class made each topic and reading seem
less intimating. Her approach in helping us to understand the
topics we discussed was very welcoming and inviting which was
very much appreciated. She would often remark on how the
readings could be dense or not easy, but that we should take out
time with them, and encouraged us to try everything.}
\item \emph{The writing workshops we had in class for our research papers
were very helpful. I also think the professor's breakdown of
essay writing was the most coherent I have ever heard. I have
never had a college course that bothered to go over the
mechanics of writing a large paper before, despite studying
English in my undergrad. The tools I have gathered from the
writing workshops and essay writings will be invaluable for me
moving forward.}
\item \emph{She made the environment of the classroom so welcoming and
safe, which I think encouraged a lot of people to be able to
speak up. She would work really hard to be inclusive, ask
questions to make us consider other POVs and also took feedback
from us to adjust assignments or due dates based on our needs.}
\end{itemize}
\end{itemize}
\end{itemize}
\section*{Mentoring \& Advising}
\label{sec:org01a31e2}
In addition to the 30 or so students that I teach directly each
semester, I currently advise 35 students, mostly in Library and
Information Science (LIS), but also in Data Analytics and
Visualization (DAV), and Information Experience Design (IXD). Each
semester, I meet individually with these students to discuss their
progress in coursework and projects, to strategize their final
portfolios, and career expectations and goals.

I also advise two students groups: the \emph{Stacks} Student Journal and
Latines in Tech. The \emph{Stacks} Student Journal is a new group founded
by students from my INFO 601: Foundations of Information course. These
students are in the process of producing their first issue of the
\emph{Stacks} journal, which will contain peer-reviewed versions of their
601 papers, showcasing the high level of critical work and writing
accomplished in that course. The second group, Latines in Tech, is
professional development group primarily focused on Latinx students.
Their main goals are to provide professional development and
networking opportunities for their members. Currently, I am working
with that group to organize a community altar for to celebrate Dia De
Los Muertos at the end of October.

\section*{Future Plans}
\label{sec:org05eedcd}
My future plans involve developing the following:
\begin{itemize}
\item Most immediately, I hope to develop, in collaboration with my
colleague Kathy Carbone, a study abroad program based in Lisbon,
Portugal. The program will likely focus on an intersection of DH and
archival work, though we are still in the early stages of discussing
the possibilities.
\item I also plan to develop a ``Readings in Queer Studies'' seminar style
course. While I am aware that theoretically heavy courses have been
less popular, I think students' appetite for theoretical material is
growing. In courses like INFO: 601 Foundations of Information and
INFO: 657 Digital Humanities, the assigned readings in Queer Studies
generate much interest and discussion, especially among students in
LIS, who often claim them as their favorite readings. For this
course, I would complement the seminar focus with fieldwork activity
that involves visiting queer archives in New York City.
\item Finally, I plan to develop a Python programming course that is
specific to working with text-based data. This course would build
from my workshop that I run for the Data Storytelling Lab, which
focuses on using advanced tools for text cleaning, processing,
annotating, and NLP analysis. This course would join my other
course, INFO 697: How to Build a Bot, to offer instruction in
intermediate programming skills.
\end{itemize}
\end{document}
